\section{Actual raw blocks and the original point lower bound}
\label{S:actual-raw-block-lower}

Fix a nonzero covariance from Hypothesis~\ref{H:homog} and any real
nonnegative nonzero $\Psi\in C_c^\infty((0,1)\times\R^d)$.
Retain the \emph{original} $Z=W(\Psi)$, $v=v_\Psi>0$ and mean $h_s$
of Lemma~\ref{L:compact-raw-profile}. Choose one finite even mask
$m_R\in[0,1]$, equal to one on $B_R$ and supported in $B_{2R}$,
as in Proposition~\ref{P:compact-original-common-trial}.
Then $V_s=m_Rh_s$ has the ground family, rate $\lambda_s$ and full
$F$ projection of Section~\ref{S:compact-mean-family}. Set
\[
  I_R=\int_0^1\lambda_s\ud s,\qquad
  Q_2=\liminf_{T\to\infty}T^{-\kappa}\log\E|u(T,0)|^2.
\]
We will prove $Q_2\ge2I_R-v/2$, for this same control and covariance.
Fix a real nonnegative smooth compact $\phi$ of integral one, the
bounded observable $O_\phi$, and the actual data $e_0=(0,1,0)\in F$.
Both displacement and velocity are retained.

\subsection{The common raw noise and its exact Gaussian shift}

Choose an integer $B\ge1$, later fixed by the cone estimate, and take
$0<\delta<\min(1,1/B)$. Define
\begin{equation}\label{E:raw-block-clocks}
  N=\left\lfloor\frac1{\delta B}\right\rfloor\ge1,\quad
  s_m=m\delta B\ (0\le m\le N),\quad
  \tau=\delta^{-1}-NB\in[0,B),\quad
  g_\delta=\delta^{(1-a)/2}.
\end{equation}
For a smooth compact positive-fast-time test $f$, extended by zero,
put
\begin{equation}\label{E:raw-block-pullback}
  T_{\delta,m}f(s,x)
  =\delta^{a/2-1}f((s-s_m)/\delta,x),\qquad
  W_{\delta,m}(f)=W(T_{\delta,m}f).
\end{equation}
Write $W_{\mathrm{fast}}=W_{\delta,0}$.
Use the same admissible $\gamma,\theta,\beta$ as
\eqref{E:driver-parameters}. Write $A_R(W_{\delta,m})$ for the
actual $F$ first driver obtained by the bounded extension
\eqref{E:constant-phase-driver}, on the local horizon $[0,B]$, and set
\begin{equation}\label{E:raw-block-drivers}
  X_m=g_\delta A_R(W_{\delta,m}),\qquad
  H_m(u)=\int_0^uS_F(-r)iB_F(V_{s_m+\delta r})S_F(r)\ud r.
\end{equation}
The convolution formula for $h$ extends smoothly to $s\in[0,2]$,
covering these means. Ground templates will only use $s_m\in[0,1]$.

\begin{lemma}[Actual coupled Gaussian block inputs]\label{L:raw-block-inputs}
  Each $T_{\delta,m}$ is an isometry of the raw test completion,
  respecting covariance-null tests. All $W_{\delta,m}$ and $Z$
  are on the same given noise space, without an independence assertion.
  Let $\mathcal K(F)$ have the operator norm and let
  $\mathscr B_0=h_0^\theta([0,B];\mathcal K(F))$ be the
  $C^\theta$ closure of zero-initial polygonal compact-operator paths.
  It is a separable Banach space. In
  $\mathscr B_\delta=\mathscr B_0^{N+1}$ with maximum norm,
  $(X,Z)$, $X=(X_m)_{m=0}^N$, is jointly centered Gaussian and
  \begin{equation}\label{E:raw-block-covariance}
    b:=\E[XZ]=\sqrt\delta\,H,\qquad H=(H_m)_{m=0}^N.
  \end{equation}
  Moreover,
  \begin{equation}\label{E:raw-block-moments}
    \|X_m\|_{L^p\mathscr B_0}
    \le C_{R,B}\sqrt p\,g_\delta,\qquad p\ge1.
  \end{equation}
\end{lemma}

\begin{proof}
  Substituting $s=s_m+\delta u$, $t=s_m+\delta r$ in the raw
  covariance gives the factor
  $\delta^{a-2}\delta^2\delta^{-a}=1$.
  Thus the maps extend isometrically through the null quotient.
  Direct test substitution gives
  $W_{\delta,m}(f)=W_{\mathrm{fast}}(f(\cdot-mB,\cdot))$.
  These are literal common-space identities.

  Finite-rank rational matrices in a countable real Hilbert basis
  are dense in $\mathcal K(F)$: for compact $T$,
  $P_nTP_n\to T$ in norm by compactness of $T,T^*$, and finite
  entries admit rational approximation. Polygonal paths with rational
  breakpoints and such matrices give separability of $\mathscr B_0$.
  To approximate breakpoints, use uniform convergence, uniform
  Lipschitz bounds and
  $[f]_\theta\le(2\|f\|_\infty)^{1-\theta}\operatorname{Lip}(f)^\theta$.
  The same inequality puts every zero-initial Lipschitz
  compact-valued path in $\mathscr B_0$ by polygonal approximation.
  The finite raw coefficient paths $iA_jL_\psi$ have these properties.
  Their actual Gaussian sums converge in $L^pC^\theta$ and hence
  in the closed space $\mathscr B_0$. Their laws and the parent
  Gaussian norm bound give~\eqref{E:raw-block-moments}.
  Every finite joint coefficient list with $Z$ is Gaussian on the
  same $W$. Scalar continuous functionals pass through the $L^2$
  limits, preserving Gaussianity by convergence of variances and
  characteristic functions. Thus $(X,Z)$ is the asserted Banach
  Gaussian variable. Cauchy--Schwarz gives
  $\E\|XZ\|\le\|X\|_2\sqrt v<\infty$, so its Bochner covariance exists.

  The original raw covariance and the test pullback give
  \begin{equation}\label{E:raw-block-mean-jacobian}
    \E[W_{\delta,m}(f)Z]
    =\delta^{a/2}\int f(u,x)h_{s_m+\delta u}(x)\ud x\ud u.
  \end{equation}
  Absolute Fubini uses compact supports, local Radon mass and $a<1$.
  This identity extends to the literal driver matrix tests on $E$:
  their time indicator is approximated by smooth bounded tests on one
  compact support, converging in the raw norm by dominated convergence
  against $|u-r|^{-a}\Gamma$; their mean pairings converge as well.
  The actual Sobolev adjoint is the one in
  \eqref{E:driver-covariance-row}. Hence the path covariance is
  $g_\delta\delta^{a/2}H_m=\sqrt\delta\,H_m$ on real Schwartz
  $E$ pairs. For constant inputs the two columns are the same
  $L_\psi$ inputs $(\psi,0)$ and $(0,\psi)$, and total finite
  propagation identifies their deterministic response.
  The span of $E$ and both constant coordinates is dense in $F$;
  $E$ alone is not. Boundedness, rational times and path continuity
  prove~\eqref{E:raw-block-covariance} on all $F$.
  In particular $H=b/\sqrt\delta\in\mathscr B_\delta$.
  This proves Lemma~\ref{L:raw-block-inputs}.
\end{proof}

\begin{proposition}[Change of law on the whole norm-Borel space]
  \label{P:raw-block-gaussian-law}
  For every nonnegative norm-Borel $j$ on all of $\mathscr B_\delta$,
  \begin{equation}\label{E:raw-block-gaussian-law}
    \E j(X)=
    \E\left[e^{-Z/\sqrt\delta-v/(2\delta)}j(X+H)\right].
  \end{equation}
  The equality is valid also as an extended nonnegative integral.
\end{proposition}

\begin{proof}
  For real $s$, the density $e^{-sZ-s^2v/2}$ defines an equivalent
  probability measure $Q_s$ by completing the square in the
  $N(0,v)$ density. For any finite list $\ell$ of continuous real
  coordinates on $\mathscr B_\delta$, the Gaussian vector
  $(\ell(X),Z)$ is equal in law to $(TG,z\cdot G)$ for a finite
  standard Gaussian vector $G$. A positive covariance square root
  gives this representation even when singular; no inverse is used.
  Here $Tz=\ell(b)$ and $|z|^2=v$. Substitution $G'=G+sz$ in its
  density gives
  $\E_{Q_s}f(\ell(X+sb))=\E f(\ell(X))$ for every bounded Borel $f$.

  The matrix entries of each block in a real Hilbert basis at
  rational times, together with $0,B$, generate the norm-Borel
  sigma-field. To see the operator-norm point, for a countable dense
  set $\mathcal D$ of unit vectors obtained from normalized rational
  finite basis combinations,
  $\|T\|=\sup_{u,w\in\mathcal D}|\langle Tu,w\rangle|$.
  These pairings are finite linear combinations of matrix entries.
  Supremum and H\"older norms of continuous paths are suprema over
  rational times and pairs. Thus distance to every fixed element
  of $\mathscr B_\delta$ is coordinate-measurable. Open balls are
  measurable by rational closed-ball approximation, and separability
  makes every open set a countable union of balls from dense centers.
  Conversely the coordinates are continuous.
  The cylinder sets form a generating pi-system; uniqueness of
  probability measures extends the finite-coordinate equality to
  all norm-Borel sets. Simple functions and monotone convergence
  give the identity for nonnegative $j$. Finally take
  $s=\delta^{-1/2}$ and use $sb=H$ from
  \eqref{E:raw-block-covariance}. The map and its deterministic
  translate are defined on the entire input space; there is no
  translated exceptional-domain assumption.
  This proves Proposition~\ref{P:raw-block-gaussian-law}.
\end{proof}

\subsection{Small noisy blocks and an indexed full-phase cone}

For any $x\in\mathscr B_0$ define
$D_x(u)=S_F(u)U_x(u,0)$ by the everywhere-defined linear Young
map of Proposition~\ref{P:constant-young-linear}.
Write $D_m^0=D_{X_m}$ and $D_m^+=D_{X_m+H_m}$.
The latter retains both the mean and the actual residual noise.

\begin{lemma}[Uniform block error]
  \label{L:raw-block-uniform-error}
  Let $M_\delta=\max_{m\le N}\|X_m\|_{\mathscr B_0}$, including
  its factor $g_\delta$. Then $M_\delta\to0$ in every finite $L^q$,
  and, for some $c>0$,
  \begin{equation}\label{E:raw-block-maximum-exponential}
    \sup_\delta\E e^{cM_\delta^2}<\infty.
  \end{equation}
  With the actual deterministic $P_s$ from the preceding section,
  \begin{equation}\label{E:raw-block-operator-error}
    E_{\mathrm{op},\delta}:=
    \max_{m\le N}\sup_{0\le u\le B}
    \|D_m^+(u)-P_{s_m}(u)\|_{\mathcal B(F)}
    \longrightarrow0
    \quad\hbox{in every finite }L^q.
  \end{equation}
\end{lemma}

\begin{proof}
  The maximum and~\eqref{E:raw-block-moments}, with $N+1\le2/\delta$,
  give
  \[
    \|M_\delta\|_p
    \le (N+1)^{1/p}C\sqrt p\,g_\delta
    \le C'\sqrt p\,\delta^{(1-a)/2-1/p}.
  \]
  For fixed $q$, choose $p>\max(q,2/(1-a))$ to prove the first limit.
  No independence is needed. Choose an integer $n_0$ with
  $n_0(1-a)\ge1$. For $n\ge n_0$,
  \[
    \E M_\delta^{2n}
    \le2C^{2n}(2n)^n\delta^{n(1-a)-1}
    \le2C^{2n}(2n)^n.
  \]
  For $1\le n<n_0$, Lyapunov bounds the moment by the appropriate
  power of the $2n_0$ moment. These finitely many terms are bounded.
  For the remaining exponential series, $n!\ge(n/e)^n$ gives
  a geometric tail $2\sum_{n\ge n_0}(2ecC^2)^n$ for small $c$.
  This proves~\eqref{E:raw-block-maximum-exponential}; no such
  assertion is made for the unscaled maximum.

  Let $H_m^{\mathrm{stat}}$ be the primitive for constant potential
  $V_{s_m}$. Compact mean smoothness on $[0,2]$ and the statewise
  primitive bound give
  \[
    \sup_m\max(\|H_m\|_{C^\theta},
    \|H_m^{\mathrm{stat}}\|_{C^\theta})\le H_0,\qquad
    \sup_m\|H_m-H_m^{\mathrm{stat}}\|_{C^\theta}\le C\delta.
  \]
  Indeed, the potential difference is at most $C\delta B$ in
  $L^\infty$ and $L^2$; integration of the bounded free groups
  gives its Lipschitz primitive bound on this fixed horizon.
  On $\{M_\delta\le1\}$, local Young stability on the fixed ball
  of radius $H_0+1$ gives
  $E_{\mathrm{op},\delta}\le C(M_\delta+\delta)$.
  This proves convergence in probability. The global Young bound
  also gives
  \[
    \max_m\sup_{u\le B}\|D_m^+(u)\|
    \le C_B\exp\big(C_B(1+H_0+M_\delta)^{1/\theta}\big).
  \]
  Since $1/\theta<2$, every finite power of this bound is dominated
  by a constant times the exponential in
  \eqref{E:raw-block-maximum-exponential}.
  Uniform higher moments yield uniform integrability, proving
  \eqref{E:raw-block-operator-error}. Suprema are measurable:
  on a countable dense set of unit states each difference is
  continuous in $u$, so rational times and those states suffice.
  Operator-norm continuity of the physical free group is not used.
  Only $s_m\le1$ supplies the ground templates.
  This proves Lemma~\ref{L:raw-block-uniform-error}.
\end{proof}

\begin{lemma}[Indexed cone and terminal displacement]
  \label{L:raw-block-indexed-cone}
  Fix $0<\varrho<1/2$. There are an integer $B\ge1$ and positive
  $\omega,\epsilon_{\mathrm{op}},c_*$ such that the following holds
  for every $N\ge1$. If $\sigma_0,\ldots,\sigma_N\in I$ have
  adjacent distances at most $\omega$ and real bounded operator
  families $D_m(u)$, $0\le u\le B$, satisfy
  \[
    \max_{m\le N}\sup_{0\le u\le B}
    \|D_m(u)-P_{\sigma_m}(u)\|\le\epsilon_{\mathrm{op}},
  \]
  then, for $0\le u\le B$,
  \begin{equation}\label{E:raw-block-positive-product}
    O_\phi D_N(u)D_{N-1}(B)\cdots D_0(B)e_0
    \ge c_*(1-\varrho)^{N-1}
    \exp\left(B\sum_{m<N}\lambda_{\sigma_m}
    +u\lambda_{\sigma_N}\right)>0.
  \end{equation}
  No positivity, composition or independence assumption on the $D_m$
  is required.
\end{lemma}

\begin{proof}
  From~\eqref{E:compact-normalized-complement} fix
  $C_{\mathrm{gap}}\ge1$, $\eta_{\mathrm{gap}}>0$ so the normalized
  complementary norm is at most
  $C_{\mathrm{gap}}e^{-\eta_{\mathrm{gap}}t}$.
  Define positive margins and finite bounds
  \[
    a_*=\inf_I\ell_s(e_0),\quad j_*=\inf_I O_\phi(e_s),\quad
    M_e=\max(1,\sup_I\|e_s\|),\quad
    M_\ell=\max(1,\sup_I\|\ell_s\|).
  \]
  Also put $M_P=\max(1,\sup_I\|\Pi_s\|)$,
  $M_Q=1+M_P$, $M_z=M_e+1$, $E_0=\|e_0\|=1$ and
  $H_{\mathrm{obs}}=\max(1,\|O_\phi\|)$.
  Choose the aperture \emph{first}:
  $\vartheta=\min(1,j_*/(16H_{\mathrm{obs}}C_{\mathrm{gap}}))>0$.
  Next choose integer $B\ge1$ with
  \[
    q_B:=C_{\mathrm{gap}}e^{-\eta_{\mathrm{gap}}B}
    \le\min\left(\frac18,\frac{\vartheta a_*}{8M_QE_0}\right).
  \]
  Choose $\nu>0$ such that
  $\nu M_z\le\min(\varrho/2,\vartheta,1/4)$.
  Uniform continuity of $\ell_s,\Pi_s$ supplies $\omega>0$
  for which adjacent distances at most $\omega$ give both norm
  differences at most $\nu$. Finally choose $\epsilon_{\mathrm{op}}>0$
  satisfying all five budgets
  \begin{align*}
    M_\ell\epsilon_{\mathrm{op}}M_z&\le\varrho/2,&
    M_Q\epsilon_{\mathrm{op}}M_z&\le\vartheta/4,\\
    M_\ell\epsilon_{\mathrm{op}}E_0&\le a_*/2,&
    M_Q\epsilon_{\mathrm{op}}E_0&\le\vartheta a_*/8,&
    H_{\mathrm{obs}}\epsilon_{\mathrm{op}}M_z&\le j_*/8.
  \end{align*}
  All denominators are positive. This order of choices uses the
  normalized complement rather than a growing terminal norm.

  The cone at $s$ consists of $z=Ae_s+r$, with
  $A=\ell_s(z)>0$, $\ell_s(r)=0$ and $\|r\|_F\le\vartheta A$.
  In particular $r=Q_sz$ and $\|z\|\le M_zA$.
  For entrance, $z_1=D_0(B)e_0$, with $b=\sigma_0$, has normalized
  amplitude at least
  $a_*-M_\ell\epsilon_{\mathrm{op}}E_0\ge a_*/2$.
  Its normalized residual is at most
  $q_BM_QE_0+M_Q\epsilon_{\mathrm{op}}E_0
  \le\vartheta a_*/4$.
  Thus it enters the cone with residual ratio at most $\vartheta/2$
  and amplitude at least $e^{\lambda_bB}a_*/2$.

  For an old-cone state $z=Ae_b+r$ and an adjacent index $c$,
  rebase before using the complementary estimate:
  \[
    \ell_c(z)\ge(1-\nu M_z)A,\qquad
    \|Q_cz\|\le\|Q_bz\|+\|(\Pi_b-\Pi_c)z\|
    \le(\vartheta+\nu M_z)A.
  \]
  For the next full block $D$ close to $P_c(B)$, put
  $x=e^{-\lambda_cB}\ell_c(Dz)$ and
  $y=e^{-\lambda_cB}\|Q_cDz\|$.
  Since $\lambda_c>0$, the perturbation and budgets give
  \[
    x\ge(1-\nu M_z-M_\ell\epsilon_{\mathrm{op}}M_z)A
    \ge(1-\varrho)A>A/2,
  \]
  and
  \[
    y\le
    [q_B(\vartheta+\nu M_z)+M_Q\epsilon_{\mathrm{op}}M_z]A
    \le\vartheta A/2.
  \]
  The last inequality uses
  $q_B\le1/8$, $0\le\nu M_z\le\vartheta$ and
  $M_Q\epsilon_{\mathrm{op}}M_z\le\vartheta/4$.
  Together with $\varrho<1/2$, these give $y\le\vartheta x$.
  These are precisely the finite residual and cone-conversion margins;
  the old residual was not assumed to lie in the new kernel.
  Each full block therefore preserves the new cone and multiplies
  its amplitude by at least $e^{\lambda_cB}(1-\varrho)$.

  For the terminal $0\le u\le B$, rebase once more to $c=\sigma_N$.
  The same bounds give $\ell_c(z)\ge3A/4$ and
  $\|Q_cz\|\le2\vartheta A$. Consequently,
  \[
    e^{-\lambda_cu}O_\phi(P_c(u)z)/A
    \ge3j_*/4-2H_{\mathrm{obs}}C_{\mathrm{gap}}\vartheta.
  \]
  The perturbation changes this by at most
  $H_{\mathrm{obs}}\epsilon_{\mathrm{op}}M_z\le j_*/8$.
  Since $16H_{\mathrm{obs}}C_{\mathrm{gap}}\vartheta\le j_*$,
  the result is at least $j_*/2$.
  Combine entrance, all $N-1$ intermediate full blocks and this
  terminal estimate to obtain~\eqref{E:raw-block-positive-product},
  with $c_*=a_*j_*/4$. The induction is empty for $N=1$ and
  the estimate includes $u=0$.
  This proves Lemma~\ref{L:raw-block-indexed-cone}.
\end{proof}

\subsection{Actual source restart and comparison with the original point}

\begin{lemma}[The actual raw restart and source products]
  \label{L:raw-block-restart}
  Fix $\delta,B,N$ and let $L=\delta^{-1}$,
  $L_{\mathrm{ext}}=(N+1)B\ge L$. On this global fast horizon define
  $A_{\mathrm{glob}}=g_\delta A_R(W_{\mathrm{fast}})$.
  Its actual physical flow obeys
  \begin{equation}\label{E:raw-block-restart}
    P_{\mathrm{glob}}(mB+u,mB)=D_m^0(u),\qquad 0\le u\le B.
  \end{equation}
  The analogous identity holds with the deterministic mean
  $V_{\delta t}$ added globally and local drivers $X_m+H_m$.
  Their finite block products, including the last remainder $\tau$,
  are actual source limits in $L^2$, at this fixed $\delta,N$.
\end{lemma}

\begin{proof}
  Take one retaining slow source $Y_{\epsilon,\eta}$ on the
  extension of the given $W$ constructed in
  Section~\ref{S:source-driver-convergence}. Define the actual
  fast and local noises by
  \[
    Y_{\mathrm{fast}}(f)=Y_{\epsilon,\eta}(T_{\delta,0}f),\qquad
    Y_m(f)=Y_{\epsilon,\eta}(T_{\delta,m}f).
  \]
  The pullback isometry and~\eqref{E:source-two-dominations} give
  source and residual variance bounds by $\|f\|_{H_W}^2$.
  The residual tends to zero on each transformed smooth compact test
  at fixed $\delta,m$. Joint Gaussianity is retained.
  Proposition~\ref{P:source-gaussian-transfer} therefore gives the
  actual $L^pC^\theta$ first-driver limits on $[0,L_{\mathrm{ext}}]$
  and on $[0,B]$, including both constant columns.

  Their pointfields are
  $Y_{\mathrm{fast}}(u,x)=\delta^{a/2}Y_{\epsilon,\eta}(\delta u,x)$
  and $Y_m(u,x)=Y_{\mathrm{fast}}(mB+u,x)$.
  Indeed, these follow by change of variables in their test integrals.
  The temporal law is $K_{\eta/\delta}$ because
  $\delta^aK_\eta(\delta r)=K_{\eta/\delta}(r)$;
  the spatial heat law and all joint width rules are unchanged.
  They are continuous in time and smooth in space.
  Their ordinary primitives equal their Gaussian-transfer drivers:
  their actual matrix entries are $Y_m(h_R)$, by the indicator
  approximation used in~\eqref{E:raw-block-mean-jacobian} and the
  point-test integral identity. Total finite propagation extends
  this identity to both constant columns.
  This does not assume a separately chosen retaining extension
  automatically commutes with time dilation.

  Changing the ordinary source integration variable and composing
  the free groups gives the exact identity on $F$
  \[
    S_F(mB)\big[A_{\mathrm{glob}}^{\mathrm{src}}(mB+u)
    -A_{\mathrm{glob}}^{\mathrm{src}}(mB)\big]
    S_F(-mB)=X_m^{\mathrm{src}}(u).
  \]
  No physical mask is commuted through a wave group.
  Restriction, difference and conjugation are bounded maps on
  these fixed-horizon $C^\theta$ spaces. Passing the actual source
  limits gives the same identity for $A_{\mathrm{glob}},X_m$.
  For the finitely many $m$, a convergent subsequence or rational
  times and path continuity gives simultaneous identities in $u$.
  Conjugation norms need only be finite at this fixed $\delta$.

  Conjugate the global interaction Young equation on
  $[mB,mB+u]$ by $S_F(mB)$. Bounded linearity of its Riemann sums
  and that raw identity give the local equation driven by $X_m$.
  Uniqueness identifies it with $U_{X_m}$, proving
  \eqref{E:raw-block-restart}. The ordinary deterministic primitive
  for $V_{\delta t}$ has the identical cocycle with local $H_m$.
  Composition now identifies the full product, with terminal $\tau$.

  For completeness, source convergence of these finite products
  requires no growing-product stability assertion. At fixed $\delta,N$
  let $C\ge1$ be the maximum norm of their source and raw factors,
  and let $e_{\mathrm{src}}$ be the maximum factor difference.
  The parent Young convergence gives every finite moment of $C$
  uniformly in source widths and $e_{\mathrm{src}}\to0$ in every
  finite $L^p$. If $d(n)$ is the difference of the first $n$
  ordered products, then
  \[
    d(n+1)\le C d(n)+C^n e_{\mathrm{src}},\qquad d(0)=0.
  \]
  Extend both finite sequences by identities so the recurrence
  holds for all natural $n$. The finite error-accumulation estimate
  gives $d(N+1)\le(N+1)C^{N+1}e_{\mathrm{src}}$.
  H\"older, using $C$ in $L^{4(N+1)}$ and the error in $L^4$,
  proves $L^2$ convergence. This is only a fixed-$\delta$ source limit.
  Shifted source products may use the fixed raw mean plus the
  regularized noise; no finite-source Cameron mean is substituted.
  This proves Lemma~\ref{L:raw-block-restart}.
\end{proof}

\begin{proposition}[The square observable and canonical comparison]
  \label{P:raw-block-point-comparison}
  Define the signed noise-only and shifted averages
  \[
    U_{\delta,B}
    =O_\phi D_N^0(\tau)D_{N-1}^0(B)\cdots D_0^0(B)e_0,
    \qquad
    Y_{\delta,B}
    =O_\phi D_N^+(\tau)D_{N-1}^+(B)\cdots D_0^+(B)e_0.
  \]
  They have all finite unweighted moments, and
  \begin{align}
    \E U_{\delta,B}^2
    &=\E\left[e^{-Z/\sqrt\delta-v/(2\delta)}
    Y_{\delta,B}^2\right],
    \label{E:raw-block-square-identity}\\
    \E U_{\delta,B}^2
    &\le\E|V_\delta(1,0)|^2
    =\E|u(T,0)|^2,\qquad \delta=T^{-\kappa}.
    \label{E:raw-block-original-comparison}
  \end{align}
  The variance $v$ in the density is the original slow raw variance.
\end{proposition}

\begin{proof}
  On all of $\mathscr B_\delta$ put
  \[
    j_\delta(x)=
    |O_\phi D_{x_N}(\tau)D_{x_{N-1}}(B)\cdots D_{x_0}(B)e_0|^2.
  \]
  The everywhere-defined continuous Young map, fixed-time bounded
  free groups, finite products and bounded $O_\phi$ make $j_\delta$
  continuous and nonnegative on the whole input space.
  Proposition~\ref{P:raw-block-gaussian-law} proves
  \eqref{E:raw-block-square-identity}. Finite products of the local
  Young flows have every finite moment by their Gaussian norm bounds
  and H\"older, at fixed $\delta,N$. The Gaussian identity itself
  did not require prior integrability of $j_\delta$.

  To apply the masked comparison at the actual amplitude, put
  $\Gamma_\delta=g_\delta^2\Gamma$ and
  $W'=g_\delta W_{\mathrm{fast}}$.
  This is a nonzero member of the same admissible covariance class,
  with unchanged degree $\alpha$ and temporal parameter $a$.
  Its actual first driver is $g_\delta$ times the original one:
  the raw matrix forms are linear in noise, and basis-independent
  Gaussian transfer and the bounded $F$ extension retain both
  constant columns. The $n$th source chain is multiplied by
  $g_\delta^n$. Thus this is exactly the noise-only model at that
  amplitude.

  Lemma~\ref{L:raw-block-restart} and the ordinary source identity
  identify $U_{\delta,B}$ as the $L^2$ limit of the tested masked
  ordinary solution driven by $W'$ at fast time $L=1/\delta$.
  Its sources $g_\delta Y_{\mathrm{fast}}$ have the correct law,
  contraction and raw test limit. Apply
  Proposition~\ref{P:masked-canonical-transfer}, or equivalently
  \eqref{E:constant-original-point-transfer}, to $\Gamma_\delta$
  and this actual source limit. Since $\int\phi=1$,
  \[
    \E U_{\delta,B}^2
    \le\E|u_{\Gamma_\delta}(L,0)|^2.
  \]
  This application describes the same noise-only fast law; it
  changes neither the original slow $\Psi,Z$ nor $v$ in the density.

  The ordinary change from fast $u$ to slow $t=\delta u$, with
  phase $(V,\delta V_t)$, identifies this unmasked fast law with
  \eqref{E:exact-fast-equation} at slow time one. Indeed,
  $g_\delta Y_{\mathrm{fast}}(u)
  =\sqrt\delta\,Y_{\epsilon,\eta}(\delta u)$,
  and its fast cutoff is $\eta/\delta$.
  The canonical fixed-order and absolutely convergent series
  construction passes this law identity at fixed $\delta$ to the
  raw moment, preserving every mixed and self-pairing term.
  Lemma~\ref{L:exact-canonical-scaling} gives the final equality
  in~\eqref{E:raw-block-original-comparison}.
  The physical scale $r_T=T\delta$ is distinct from the fixed
  mask radius $R$. The mask need not contain the growing backward
  cone: the comparison is a moment inequality, not equality of
  masked and full fields. All source limits preceded this large-time
  comparison.
  This proves Proposition~\ref{P:raw-block-point-comparison}.
\end{proof}

\begin{proposition}[Original point lower rate at the raw half-cost]
  \label{P:raw-block-original-lower}
  For every original positive $\Psi$ and eligible fixed mask $R$
  chosen above, and every fixed $x\in\R^d$,
  \begin{equation}\label{E:raw-block-original-lower}
    \liminf_{T\to\infty}T^{-\kappa}\log\E|u(T,x)|^2
    \ge 2\int_0^1\lambda_s\ud s-\frac{v_\Psi}{2}.
  \end{equation}
\end{proposition}

\begin{proof}
  Fix $e_{\mathrm{rate}}>0$. Choose $0<\varrho<1/2$ with
  $\log(1-\varrho)>-e_{\mathrm{rate}}/3$, then choose the cone
  constants, including fixed $B$, from
  Lemma~\ref{L:raw-block-indexed-cone}.
  For small $\delta$ the indices $s_m=m\delta B$ lie in $I$
  with adjacent gaps at most $\omega$.
  Define the measurable common-space event
  \[
    E_\delta=\{E_{\mathrm{op},\delta}\le\epsilon_{\mathrm{op}}\}
    \cap\{|Z|\le\delta^{-1/4}\}.
  \]
  Lemma~\ref{L:raw-block-uniform-error} and Chebyshev give
  $\mathbb P(E_\delta)\to1$, because
  $\mathbb P(|Z|>\delta^{-1/4})\le v\delta^{1/2}$.
  This uses a union bound, with no independence from $Z$ or
  between blocks.

  On $E_\delta$ the shifted blocks meet the cone hypotheses.
  In particular $Y_{\delta,B}>0$, and
  \[
    \delta\log Y_{\delta,B}
    \ge\delta\log c_*+\delta(N-1)\log(1-\varrho)
    +\delta B\sum_{m<N}\lambda_{s_m}
    +\delta\tau\lambda_{s_N}.
  \]
  The last two terms form the left Riemann sum for $I_R$ on
  the partition including its final interval of length
  $\delta\tau$, with mesh at most $\delta B$.
  The actual rate is Lipschitz, so its error is at most $C_R\delta B$.
  Also $0\le\delta(N-1)\le1/B\le1$ and hence the logarithmic
  loss is at least $\log(1-\varrho)>-e_{\mathrm{rate}}/3$.
  For sufficiently small $\delta$, the Riemann error and
  $\delta|\log c_*|$ are each at most $e_{\mathrm{rate}}/3$.
  Thus
  $Y_{\delta,B}\ge\exp((I_R-e_{\mathrm{rate}})/\delta)$ on $E_\delta$,
  including the $N=1$ and $\tau=0$ cases of the cone estimate.

  The square integrand in~\eqref{E:raw-block-square-identity}
  is nonnegative. Restrict it to $E_\delta$, where
  $-Z/\sqrt\delta\ge-\delta^{-3/4}$.
  With $p_\delta=\mathbb P(E_\delta)$ this gives
  \[
    \E U_{\delta,B}^2
    \ge p_\delta
    \exp\left(\frac{2I_R-2e_{\mathrm{rate}}-v/2}{\delta}
    -\delta^{-3/4}\right).
  \]
  For small $\delta$, $p_\delta\ge1/2$ and
  $\delta\log p_\delta\to0$. The normalized moderate penalty is
  $\delta(\delta^{-1/4}/\sqrt\delta)=\delta^{1/4}\to0$.
  Taking logarithms and lower limits at this fixed
  $e_{\mathrm{rate}}$ and its associated $B$ gives
  \[
    \liminf_{\delta\downarrow0}\delta\log\E U_{\delta,B}^2
    \ge2I_R-2e_{\mathrm{rate}}-v/2.
  \]
  By~\eqref{E:raw-block-original-comparison}, the
  $B$-independent original point rate is at least this right side
  for every fixed $e_{\mathrm{rate}}$. Only now let
  $e_{\mathrm{rate}}\downarrow0$. The bijection
  $\delta=T^{-\kappa}$ covers all sufficiently large $T$.

  Finally, spatial translation of all positive-width vertices
  preserves the original noise law, Green kernels and constant
  initial data. The fixed-order $L^2$ and absolute-series limits
  preserve equality of point second moments. Hence zero can be
  replaced by every fixed $x$ in~\eqref{E:raw-block-original-lower}.
  Every control and mask were fixed before asymptotics; a supremum
  over them may follow their individual inequalities.
  This proves Proposition~\ref{P:raw-block-original-lower}.
\end{proof}
